\section{Definite pairs and the coefficient sphere}\label{sec:definite-pairs}

We first associate an almost-complex structure to a definite pair and
identify the condition on a third form that makes the triple positive.
This converts the later deformation problem into the choice of taming
forms in a prescribed class. We also establish the topological
restrictions and the degree-one sphere of structures needed to produce
curves in \Cref{sec:families}.

\begin{definition}
A pair $(A,C)$ of real two-forms is \emph{definite} if its wedge-product
matrix is positive definite at every point. A two-form $B$ \emph{tames}
an almost-complex structure $J$ if $B(v,Jv)>0$ for every nonzero tangent
vector $v$. It is \emph{compatible} if, in addition,
$B(Ju,Jv)=B(u,v)$. When these forms are called symplectic, closedness is
included.
\end{definition}

\begin{lemma}\label{lem:definite-pair}
A definite pair $(A,C)$ on an oriented four-manifold determines an
almost-complex structure up to sign. Its anti-invariant real two-form
bundle is pointwise spanned by $A,C$, and its canonical complex line
bundle is trivial. A third form $B$ completes a positive triple if and
only if it tames one of the two signs of this almost-complex structure.
On a connected manifold the sign is fixed by $B$.

A real two-plane is Lagrangian for both $A$ and $C$ if and only if it is
a complex line, without specifying its orientation.
\end{lemma}

\begin{proof}
Ratios of positive four-forms are smooth functions. Set
\begin{equation}\label{eq:pair-normalization}
 x=\frac{AC}{A^2},\qquad
 y=\left(\frac{C^2}{A^2}-x^2\right)^{1/2}>0,\qquad
 T=\frac{C-xA}{y}.
\end{equation}
Then $AT=0$ and $T^2=A^2>0$. Consequently
\[
 \Omega=A+iT,\qquad \Omega^2=0,\qquad
 \Omega\overline\Omega=2A^2>0.
\]
The form $\Omega$ is complex decomposable and its kernel in the
complexified tangent bundle is complementary to its conjugate. Taking
this kernel as $T^{0,1}$ defines an almost-complex structure with the
given orientation. Replacing $y$ by $-y$ conjugates $\Omega$ and changes
$J$ to $-J$. The real and imaginary parts of $\Omega$ span the
anti-invariant forms, and $\Omega$ is a nowhere-zero section of the
canonical bundle.

Choose any Hermitian metric for $J$. The invariant two-forms are the
wedge-orthogonal complement of $\operatorname{span}(A,C)$ and have
signature $(1,3)$. For the decomposition
$B=B^{1,1}+B^{2,0+0,2}$, the positive-triple condition is exactly
$(B^{1,1})^2>0$. A real $(1,1)$ form in complex dimension two has positive
square precisely when its two Hermitian eigenvalues have the same
strict sign. The positive component is the cone of forms positive on
every complex line. The anti-invariant summand vanishes on such lines,
so this is equivalent to taming by $B$. The other component is positive
for $-J$. Connectedness fixes the choice continuously and globally.

Finally, for independent real vectors $u,v$,
$A(u,v)=C(u,v)=0$ is equivalent to $\Omega(u,v)=0$.
In a complex two-dimensional vector space with a nondegenerate complex
volume form, this says that $u,v$ are complex dependent, equivalently
that their real span is a complex line.
\end{proof}

In particular, for fixed $J$ the space of taming two-forms is an open
convex cone. Its intersection with the closed representatives of any
one class is also convex. The associated notion of positivity will
always use the sign of $J$ fixed by a stated tamer.

\begin{proposition}\label{prop:topology}
If $X$ carries a positive closed triple, then
\[
 b^+(X)=3,\qquad
 (b_1(X),b^-(X))=(0,19)\ \text{or}\ (4,3).
\]
The free integral intersection lattice is respectively
$3H\oplus2(-E_8)$ or $3H$, where $H$ is the hyperbolic plane.
Every almost-complex structure obtained from a definite pair in the
triple has vanishing integral first Chern class.
\end{proposition}

\begin{proof}
By \Cref{lem:definite-pair}, one member of the triple is a symplectic
form taming an almost-complex structure whose canonical bundle is
trivial. The first Chern class of this symplectic form therefore
vanishes: the space of its tamed almost-complex structures is
contractible and contains compatible structures. Bauer's bound for a
closed symplectic four-manifold with torsion first Chern class gives
$b^+\le3$ \cite[Corollary~1.2]{Bauer2008}. The three independent
positive cohomology classes give the reverse inequality.

The characteristic-number identity $c_1^2=2\chi+3\sigma$, obtained
from $p_1=c_1^2-2c_2$, the Euler-number identity and the signature
theorem \cite[Corollaries~11.12 and~15.5, Theorem~19.4]{MilnorStasheff1974},
together with $b^+=3$, gives
\[
 0=19-4b_1-b^-,\qquad \sigma=4b_1-16.
\]
Moreover $w_2\equiv c_1=0\pmod2$, so $X$ is spin
\cite[Problem~14-B]{MilnorStasheff1974}. Rokhlin's theorem makes
$\sigma$ divisible by $16$ \cite{Rokhlin1952}.
Since $b^-\ge0$, these relations give
$b_1\in\{0,4\}$ and the displayed alternatives. Poincar\'e duality
makes the free intersection form unimodular, and the Wu formula makes
it even \cite[Theorem~11.14]{MilnorStasheff1974}. Classification of
indefinite even unimodular lattices now gives the two stated forms
\cite[Chapter~V, Section~2.2, Theorems~5--6]{Serre1973}.
Only the intersection lattice has been classified at this stage.
\end{proof}

The positive three-plane bundle spanned by the triple is the
self-dual bundle of a unique oriented conformal structure. Choose one
metric $g_0$ in that conformal class. Since the original forms are
closed and self-dual, they are harmonic, and
\begin{equation}\label{eq:pair-harmonic-basis}
 \mathcal H^+_{g_0}
 =\operatorname{span}_{\R}\{\omega_1,\omega_2,\omega_3\}
\end{equation}
by \Cref{prop:topology}. Here the span is the vector space of global
forms, with constant coefficients.

For $v\in S^2\subset\R^3$, let $B_v=v\cdot\omega$ and let
$\mathcal A_v$ be the plane of forms $a\cdot\omega$ with $a\perp v$.
Choosing the sign tamed by $B_v$ defines an almost-complex structure
$J_v$. This construction uses the plane $\mathcal A_v$ itself and does
not require a global choice of its ordered basis over $S^2$.

\begin{lemma}\label{lem:twistor-degree}
The family $J_v$ is smooth and is Hermitian for $g_0$. At every point
of $X$ its map to the twistor sphere has degree of absolute value one.
The closed $J_v$-anti-invariant forms form the two-dimensional space
$\mathcal A_v$.
\end{lemma}

\begin{proof}
The first assertion follows from the construction and the
self-dual description of Hermitian almost-complex structures in
dimension four. At a fixed point, let $K$ be the positive Gram matrix
of the triple for $g_0$. A vector normal to the coefficient plane
$v^\perp$, for this Gram matrix, has coefficients $K^{-1}v$.
Its pairing with $B_v$ is positive. Thus, after normalization, the
fundamental forms give the map
\[
 v\longmapsto\frac{K^{-1}v}{|K^{-1}v|}
\]
in radial coefficient coordinates. Positive definite matrices connect to the
identity through positive definite matrices, so this map has degree
one. The orientation convention identifying coefficient and twistor
spheres can change only its sign.

A closed anti-invariant form is self-dual for $g_0$, hence harmonic.
By \eqref{eq:pair-harmonic-basis} it is a constant combination of the
original triple. It is anti-invariant for $J_v$ exactly when that
coefficient vector lies in $v^\perp$. This proves the final assertion.
\end{proof}

After a constant orthogonal change of coefficients we will denote a
chosen original ordered triple by $(A,C,B)$. For all later exact
deformations of the pair, set
\begin{equation}\label{eq:pair-lorentz-space}
 V=[A]^\perp\cap[C]^\perp\subset H^2(X;\R).
\end{equation}
This space has signature $(1,b^-)$, and $[B]\in V$ is a unit timelike
vector. It determines the future component of the timelike cone.

\begin{lemma}\label{lem:complex-symplectic}
Let $A,D$ be closed real two-forms with
$AD=0$ and $A^2=D^2>0$ pointwise. Then $A+iD$ defines an integrable
complex structure, and is a nowhere-zero holomorphic two-form for it.
If a closed third form $C$ has $AC=DC=0$ and $C^2=A^2$, then the
resulting triple is hyperk\"ahler, after choosing the sign of the
complex structure so that $C$ is positive.
\end{lemma}

\begin{proof}
Set $\Omega=A+iD$ and use its kernel as $T^{0,1}$ as in
\Cref{lem:definite-pair}. For vector fields $U,V$ in that kernel,
Cartan's formula and $\dd\Omega=0$ give
\[
 \mathcal L_U\Omega=0,\qquad
 \iota_{[U,V]}\Omega
 =\mathcal L_U(\iota_V\Omega)-\iota_V(\mathcal L_U\Omega)=0.
\]
The kernel is involutive. The smooth Newlander--Nirenberg theorem
therefore gives integrability \cite{NewlanderNirenberg1957}, and
closedness of the $(2,0)$ form
$\Omega$ gives holomorphicity.

Under the additional assumptions, $C$ is a positive closed $(1,1)$
form, hence a K\"ahler form. Its metric volume is $C^2/2=A^2/2$.
The ratio $\Omega\overline\Omega/C^2$ is constant, so the holomorphic
volume form has constant pointwise norm. The Chern connection of the
canonical line then makes $\Omega$ parallel. For a K\"ahler metric
this is the connection induced by the Levi-Civita connection.
Thus $A,D,C$ are parallel; their normalized self-dual algebra gives
the hyperk\"ahler structure. Explicitly, if
$(\eta_1,\eta_2,\eta_3)=(A,D,C)$ and $\mu=A^2/2$, then
$\eta_i\wedge\eta_j=2\delta_{ij}\mu$; thus their normalized
matrix $\mathcal Q$ from the introduction is the identity.
\end{proof}
